\documentclass[12pt,a4paper]{extarticle}
\usepackage{pdfpages}
\usepackage[utf8]{inputenc}
\usepackage{multicol,multirow}
 \usepackage{charter}
\usepackage[
a4paper,
twoside,
bindingoffset=0.5cm,
inner=2cm,
outer=2cm,
top=3cm,
bottom=3cm,
headsep=1.2cm
]{geometry}
\usepackage[sort,numbers]{natbib}
\usepackage{etoolbox}
\usepackage{bibentry}
% \usepackage{nyul_thesis}
\usepackage{graphicx}
\usepackage{caption}
\usepackage{setspace}
\usepackage{hyperref}
\usepackage{enumitem}
\usepackage{float}
\usepackage{framed}
\usepackage{titlesec} % package to reformat section titles
%\usepackage[euler]{textgreek}


\newcommand{\thesispointone}{A methodology and case study for code synthesis \\ \mbox{evaluation} of LLMs}
\newcommand{\thesispointtwo}{A program synthesis dataset for LLM temperature \\ \mbox{analysis}}
\newcommand{\thesispointthree}{Assessment of GPT-4 in Fixing Real-World Software \\ \mbox{Vulnerabilities}}

\newcommand{\thesispointoneuf}{A methodology and case study for code synthesis evaluation of LLMs}
\newcommand{\thesispointtwouf}{A program synthesis dataset for LLM temperature analysis}
\newcommand{\thesispointthreeuf}{Assessment of GPT-4 in Fixing Real-World Software Vulnerabilities}

\newcommand{\thesispointonehu}{Módszertan és esettanulmány a nagy nyelvi modellek kódszintézis-értékeléséhez}
\newcommand{\thesispointtwohu}{Programszintézis adathalmaz a nagy nyelvi modellek temperature elemzéséhez}
\newcommand{\thesispointthreehu}{A GPT-4 értékelése valós szoftversérülékenységek \\ javításában}

\newcommand{\thesisonetext}{I reviewed existing LLM comparison methodologies and identified their limitations, which led me to design a more comprehensive evaluation approach.
I selected a suitable benchmark, evaluated multiple language models on it, and performed a quality analysis of the synthesized programs.
In addition, I designed and implemented a human evaluation framework to assess developer acceptance of LLM-generated code and demonstrated the methodology through a case study.}

\newcommand{\thesistwotext}{I designed the study to create the dataset containing 18,900 raw and 18,896 processed C++ code generation outputs from large language models.
I developed a model selection framework to systematically identify suitable open-source models and evaluated nine LLMs from three model families with parameter sizes up to 70 billion using varying temperature settings.
Finally, I analyzed the results to investigate the effect of the temperature hyperparameter on source code generation and showed that it does not reveal a clear performance pattern.}


\newcommand{\thesisthreetext}{I explored prompt engineering techniques to generate fixed source code with GPT-4 for automatic program repair, aiming to maintain functionality while removing vulnerabilities.
I designed the evaluation framework for both prompts and models, using the Vul4J benchmark and an additional post-training benchmark to ensure unbiased assessment.
I also participated in the manual evaluation of the generated source code and textual responses, finding that GPT-4 could fix 33.33\% of vulnerabilities and provide useful guidance in roughly 50\% of cases.}

\pagestyle{plain}

\captionsetup[figure]{labelfont={bf},textfont={it}}
\captionsetup[table]{labelfont={bf},textfont={it}}

\begin{document}

\nobibliography*

\newcommand{\doctype}{PhD disszertáció összefoglaló}
%\include{titlepage}

%%%%%%% title page


\thispagestyle{empty}

\newcommand{\dolgozatcim}{Nagy Nyelvi Modellek Képességeinek Vizsgálata a Szoftverfejlesztésben}

\begin{center}
	%	\vfill
	\vspace*{0.25cm}
	
	\begin{spacing}{2}
		{\Huge \textbf{\dolgozatcim}}
	\end{spacing}
	
	\vspace*{2cm}
	
	{\Large \doctype}
	
	%	\vspace{2.5cm}
	\vfill
	
	{\Large Zoltán Ságodi}
	
	\vspace{0.7cm}
	
	{\Large Témavezetők: \\ \vspace{0.5em} Siket István, PhD\\ \vspace*{1em} Ferenc Rudolf, PhD}
	
	%	\vspace{2cm}
	\vfill
	
	{\large Informatika Doktori Iskola}
	
	\vspace{0.25cm}
	
	{\large Szoftverfejlesztés Tanszék}
	
	\vspace{0.25cm}
	
	{\large Természettudományi és Informatikai Kar}
	
	\vspace{0.25cm}
	
	{\large Szegedi Tudományegyetem}
	
	%	\vspace{4cm}
	\vfill
	
	\includegraphics[width=0.3\linewidth]{Figures/szte_logo}
	
	\vfill
	
	{\large Szeged
		
		\today}
	
\end{center}





%%%%%%%% title page end
 
\newpage
\thispagestyle{empty}
\mbox{}

\newpage

\setcounter{page}{1}
 
\titleformat{\section}[block]
{\normalfont\Large\bfseries}
{Tézis \thesection: }
{0pt}
{}

\renewcommand{\thesection}{\arabic{section}}

\newpage
\section*{Bevezetés}

%Software engineering is a continuously growing and evolving field.
%Not only do new subfields emerge constantly, but the underlying technologies are also changing at a rapid pace.
%From the early days of manually written software, we have progressed to a point where Integrated Development Environments (IDEs)—which provide substantial support through features such as syntax highlighting, code navigation, and automated formatting—have become standard tools in every developer’s arsenal.
%Although these tools were created for developer convenience, many have failed over time, and those that remain continue to evolve in response to the changing requirements of the developer community.
A szoftverfejlesztés folyamatosan növekvő és gyorsan változó terület.
Nemcsak új rész\-te\-rü\-le\-tek jelennek meg egyre gyakrabban, hanem az alapjául szolgáló technológiák is rendkívül gyors ütemben fejlődnek.
A kezdeti, szövegszerkesztőkben írt programoktól, eljutottunk oda, hogy az integrált fejlesztőkörnyezetek (IDE-k) – amelyek olyan funkciókkal segítik a munkát, mint a szintaxiskiemelés, a kódnavigáció vagy az automatikus formázás – mára minden fejlesztő alapvető eszközeivé váltak.
Bár ezeket az eszközöket a fejlesztők mun\-ká\-já\-nak megkönnyítésére hozták létre, sok közülük idővel eltűnt, a fennmaradók pedig folyamatosan alkalmazkodnak a fejlesztői igények változásaihoz.


%IDEs rely on a variety of analysis techniques, with source code analysis being a central component.
%While some features may appear straightforward, such as navigating from a caller to a callee, they often depend on extensive research, such as the construction of call graphs.
%Building call graphs is a complex task, as demonstrated in several of our studies~\cite{cg1, cg2, cg3}.
Az IDE-k működése számos elemzési technikára épül, amelyek közül a for\-rás\-kód-elem\-zés kiemelt szerepet játszik.
Egyes funkciók első ránézésre egyszerűnek tűnhetnek – pél\-dá\-ul a hívó és a hívott függvény közötti navigáció –, valójában azonban komoly kutatási e\-red\-mé\-nyek állnak mögöttük, például a hívási gráfok előállítása.
A hívási gráfok felépítése összetett feladat, amint azt több korábbi munkánk is bemutatja~\cite{cg1, cg2, cg3}.


%Just as a small subset of IDE features is underpinned by a substantial research foundation, newly developed tools also require a solid research base to ensure their effectiveness and reliability.
%Such tools are the recently introduced Large Language Models (LLMs), which are currently reshaping the daily work of software engineers.
%As with earlier tools, these models must adapt to developer needs, and developers, in turn, must explore their capabilities to determine which tasks can be performed at higher quality—or at the same quality more efficiently—through their use.
Ahogyan az IDE-k egyes funkciói mögött jelentős kutatási háttér áll, úgy az új eszközök esetében is elengedhetetlen a megfelelő tudományos megalapozottság a hatékony és meg\-bíz\-ha\-tó működéshez.
Ilyen új eszközök a nagy nyelvi modellek (LLM-ek), amelyek napjainkban alapvetően alakítják át a szoftverfejlesztők mindennapi munkáját.
Ezeknek a modelleknek is alkalmazkodniuk kell a fejlesztői igényekhez, miközben a fejlesztőknek fel kell fedezniük a bennük rejlő lehetőségeket annak érdekében, hogy meghatározzák, mely feladatok végezhetők el jobb minőségben – vagy azonos minőség mellett hatékonyabban – a segítségükkel.



%Developers today face numerous challenges, including the need to change established coding habits and adapt to a workflow in which creative work is increasingly transformed into review-oriented work through the use of AI models.
%This work also highlights additional challenges that are primarily technical in nature.
%With new LLMs appearing almost daily, developers struggle to select the most suitable model for their tasks.
%Experimenting with the most popular models is time-consuming and may lead them to overlook potentially optimal alternatives.
%A further issue is that newly introduced tools and models lack a common, standardized framework for quality measurement and comparison; as a result, developers cannot confidently prefer one model over another.
%Although various benchmarks are available, they mostly measure the functionality of the generated code.
%Even with versatile benchmarks, developers have to find those benchmarks and often evaluate the models on them, as most of the models are not evaluated from the perspectives of source code quality, security, or other non-functional criteria.
A fejlesztők ma számos kihívással szembesülnek.
Egyrészt szükségessé válik a meg\-szo\-kott programozási gyakorlatok átalakítása, másrészt egy olyan munkafolyamathoz kell alkalmazkodniuk, amelyben az alkotó jellegű munka egyre inkább ellenőrzésközpontú te\-vé\-keny\-ség\-gé alakul az AI-alapú eszközök használata miatt.
A dolgozat több, elsősorban technikai jellegű, problémára is rámutat.
Az egyik ilyen kihívás, hogy a szinte naponta megjelenő új LLM-ek közül nehéz kiválasztani az adott feladathoz legmegfelelőbbet.
A legnépszerűbb modellek kipróbálása időigényes, és könnyen előfordulhat, hogy így a ke\-vés\-bé ismert, de potenciálisan jobb modellek kimaradnak.
További probléma, hogy hi\-ány\-zik egy egységes, szabványosított keretrendszer a modellek minőségének mérésére és össze\-ha\-son\-lí\-tá\-sá\-ra, így a fejlesztők nem tudnak megalapozott döntést hozni egyik vagy másik modell mellett.
Bár léteznek különböző benchmarkok, ezek többnyire a generált kód funkcionalitását mérik.
Még a komplexebb benchmarkok esetében is, a fejlesztőknek kell azokat felkutatniuk, és gyakran saját maguknak kell lefuttatniuk a kiértékeléseket, mivel a modelleket ritkán vizsgálják a forráskód minősége, biztonsága vagy más nem-funkcionális szempontok alapján.


%As benchmarks are mentioned, another technical challenge is that, although some benchmarks are available, developers cannot fully trust them.
%These benchmarks often rely on small, non–real-world scenarios.
%Such benchmarks do not provide actual information about how models would perform when applied in actual development.
%It is also a problem that these benchmarks are prone to overfitting, as models may be optimized specifically for better benchmark performance.
%As model trainers know which benchmarks are checked for quality measures, they might place those benchmarks in the training set multiple times in order to seemingly increase the performance.
%Consequently, developers need reliable insight into the true real-world capabilities of LLMs.
A benchmarkokkal kapcsolatban további problémák is felmerülnek.
Bár rendelkezésre állnak, a fejlesztők nem bízhatnak meg bennük teljes mértékben.
Gyakran kis méretű, nem valós környezetből származó példákra épülnek, így nem adnak hiteles képet arról, hogyan teljesítenének a modellek valós fejlesztési helyzetekben.
Emellett ezek a benchmarkok hajlamosak a túlillesztésre: a modellek kifejezetten a jobb benchmark-eredmények elérésére optimalizálhatók.
Mivel a modellfejlesztők tisztában vannak azzal, hogy mely benchmarkokat használják kiértékelésre, előfordulhat, hogy ezeket többször is beépítik a tanítóadatokba, ezzel mesterségesen javítva az eredményeket.
Ennek következtében a fejlesztőknek megbízhatóbb képre van szükségük az LLM-ek valós, gyakorlati tel\-je\-sít\-mé\-nyé\-ről.

%An additional problem is that these models are controlled by numerous hyperparameters.
%While these offer extensive customization once a well-performing model has been selected, developers often lack guidance on how to configure the model effectively for their specific tasks or whether consistent configuration patterns exist across different models.
%Although general information about many hyperparameters is available at a broader level, everyday users frequently cannot benefit from, or may not even notice, their impact in practice.
További nehézséget jelent, hogy ezek a modellek számos hiperparaméterrel ren\-del\-kez\-nek.
Bár ezek lehetővé teszik a finomhangolást egy jól működő modell esetén, a fejlesztők gyakran nem kapnak egyértelmű útmutatást arra, hogyan érdemes ezeket beállítani konk\-rét feladatokhoz, illetve hogy léteznek-e általános, modellek között is érvényes beállítási min\-ták.
Habár általános leírások sok esetben állnak rendelkezésre, a mindennapi használat során ezek hatása gyakran nem egyértelmű vagy nehezen érzékelhető.

%While these challenges may be approached from multiple subfields of software engineering, the primary focus of this work is on source code generation, source code quality, and source code security.
%Source code generation and quality form the foundation of all software development workflows, as software cannot exist without code, and quality assurance is a routinely performed and essential activity in large-scale projects.
%Source code security is a well-established research area; however, its importance continues to increase as the frequency and impact of security incidents grow steadily over time.
Ezek a kihívások a szoftverfejlesztés több területéről is megközelíthetők, jelen dolgozat elsősorban a forráskód-generálásra, a forráskód minőségére és a forráskód biztonságára fókuszál.
A forráskód-generálás és a minőség alapvető szerepet játszik minden fejlesztési folyamatban, hiszen szoftver nem létezhet kód nélkül, és a minőségbiztosítás nagy rend\-sze\-rek esetében mindennapos, elengedhetetlen tevékenység.
A forráskód biztonsága pedig régóta kutatott terület, amelynek jelentősége folyamatosan növekszik a biztonsági incidensek számának és súlyosságának emelkedésével.

%In the following three sections, the thesis points are presented, each addressing one of the challenges outlined above.
A következő három fejezetben a tézisek bemutatása következik, amelyek mindegyike a fent ismertetett kihívások egy-egy aspektusára ad választ.


\section{\thesispointonehu}

%Large Language Models are deployed on a daily basis, therefore, developers cannot test every model in order to find the best for their usecase.
%A methodology must be applied to get a reliable and uniform comparison.
%We\footnote{The author states that while the thesis results are primarily his own work, the
%	pronoun We is used instead of I to recognize the input of co-authors in the papers
%	forming the basis of this thesis.} created a comparison methodology based on various papers that compare LLMs~\cite{copilot_code_evaluation,copilot_code_evaluation2, copilot_code_evaluation_3}.
%This way, we came-up with a methodology that takes multiple aspects into consideration not only a few.
%The methodology details how every step should be performed what must be taken into consideration.
%The methodology is also showcased in a case study to provide a working example on how to use it.
Naponta jelennek meg az újabb nagy nyelvi modellek (LLM), ezért a fejlesztők nem tudják az összes modellt egyenként kipróbálni annak érdekében, hogy megtalálják a számukra legmegfelelőbbet, hiszen nincsen rá elegendő idő és erőforrás.
Egy megbízható és egységes össze\-ha\-son\-lí\-tás érdekében szükség van egy megfelelő módszertan alkalmazására.
Kü\-lön\-bö\-ző LLM-ek össze\-ha\-son\-lí\-tá\-sá\-val foglalkozó tanulmányok alapján~\cite{copilot_code_evaluation,copilot_code_evaluation2, copilot_code_evaluation_3}, egy össze\-ha\-son\-lí\-tá\-si módszertant dolgoztunk\footnote{A szerző kijelenti, hogy bár a dolgozat eredményei elsősorban saját munkáján alapulnak, az egyes szám első személy helyett többes szám első személyt használ (''mi''), ezzel is elismerve a dolgozat alapját képező publikációk társszerzőinek hozzájárulását.} ki.
Ennek eredményeként egy olyan megközelítést hoztunk létre, amely nem csupán néhány szempontot, hanem több különböző aspektust is figyelembe vesz.
A módszertan részletesen meghatározza az egyes lépések végrehajtását, valamint azt is, hogy milyen tényezőket kell figyelembe venni.
A módszertant egy esettanulmányon keresztül is bemutatjuk, szem\-lél\-tet\-ve annak gyakorlati alkalmazását.


%The methodology starts with the right prompt selections.
%In contrast to methods everyone is used to, LLMs cannot be used without proper prompting.
%Denny et al.~\cite{codex_prompting} showed that Natural Language prompting has the effect to change the generated source code.
%The first thing we must do is to select the right prompt for the model, for the task We are going to use the model for.
%Selecting the right prompt is not easy.
%There are guidelines that should be followed, such as the work of White et al.~\cite{chatgpt_prompt_design} or the advice on GitHub~\cite{github_prompt_engineering}.
A módszertan első lépése a megfelelő prompt kiválasztása.
A hagyományos mód\-sze\-rek\-kel ellentétben az LLM-ek nem használhatók hatékonyan megfelelő prompt nélkül.
Denny és társai~\cite{codex_prompting} kimutatták, hogy a természetes nyelvű promptok jelentősen befolyásolják a generált forráskódot.
Ezért első lépésként ki kell választani a megfelelő promptot az adott modellhez és az adott feladathoz.
A megfelelő prompt kiválasztása nem triviális feladat.
Ehhez különböző irányelvek állnak rendelkezésre, például White és társai munkája~\cite{chatgpt_prompt_design}, illetve a GitHub által közzétett ajánlások~\cite{github_prompt_engineering}.


%After selecting the right prompt, we must evaluate the functional validity of the models.
%Functional validity, simply said, measures whether the model can provide code that does what it should.
%The hardest task at this point is to find a way to measure that it does ,,what it should do''.
%It is most likely to be performed via automatized tests, although it is important to note that these tests might not take the generated code as is, so a framework might be needed.
%The use of frameworks depends entirely on the requirements because they may require that the code be testable in existing tests.
A megfelelő prompt kiválasztását követően a modellek funkcionális helyességét kell értékelni.
A funkcionális helyesség azt méri, hogy a modell képes-e olyan kódot előállítani, amely valóban a kívánt módon működik.
Ennek mérése azonban nem egyszerű, mivel meg kell határozni, mit jelent pontosan az, hogy a kód ''azt csinálja, amit kell''.
Ez általában automatizált tesztek segítségével történik, ugyanakkor fontos megjegyezni, hogy ezek a tesztek nem mindig képesek közvetlenül kezelni a generált kódot, ezért adott esetben szükség lehet egy köztes keretrendszerre.
A keretrendszer alkalmazása a konkrét követelményektől függ, például attól, hogy a kódnak illeszkednie kell-e meglévő teszt\-kör\-nye\-ze\-tek\-hez.


%Getting the tests is not an easy task either, as it is either up to the team to create the tests or they must use existing benchmarks; although in that case, overfitting must be taken into consideration.
%Either way, tests are not easy to obtain in good quality.
%In our case study, to get tests that are optimal for this evaluation, we used PSB2~\cite{codegen_benchmark}, which includes 25 programming tasks.
A tesztek előállítása önmagában is kihívást jelent.
Ez történhet saját tesztek írásával, vagy meglévő benchmarkok felhasználásával.
Utóbbi esetben azonban figyelembe kell venni a túlillesztés (overfitting) lehetőségét.
Mindkét megközelítés esetén nehéz jó mi\-nő\-sé\-gű teszteket biztosítani.
Az esettanulmányunkban a PSB2 benchmarkot használtuk~\cite{codegen_benchmark}, amely 25 programozási feladatot tartalmaz, és alkalmas a modellek értékelésére.


%At the third step of the methodology, we must evaluate the non-functional validity or technical quality.
%This step ensures that the introduced LLM maintains, or rather increases, the quality of the source code.
%It can be measured, for example, by SonarQube, which we used in the case study.
A módszertan harmadik lépése a nem-funkcionális megfelelőség, azaz a technikai mi\-nő\-ség értékelése.
Ez a lépés biztosítja, hogy az alkalmazott LLM fenntartsa, vagy lehetőség szerint javítsa a forráskód minőségét.
A minőség mérhető például statikus kódelemző eszközökkel; esettanulmányunkban a SonarQube rendszert alkalmaztuk.

%As our last step, we included the human evaluation in order to see if developers are willing to accept source code generated by LLMs.
%This step, although difficult and expensive to perform, provides a better overview of whether an LLM should be introduced to the development team or not.
%In our case study, We found that, although developers typically accept the generated code, it often requires additional refinement before final use.
A módszertan utolsó lépéseként a humán értékelést is bevontuk annak érdekében, hogy felmérjük, a fejlesztők mennyire fogadják el az LLM-ek által generált kódot.
Ez a lépés ugyan költséges és nehezen kivitelezhető, ugyanakkor átfogóbb képet ad arról, hogy érdemes-e egy adott modellt bevezetni a fejlesztési folyamatba.
Eredményeink alapján a fejlesztők általában elfogadják a generált kódot, azonban az gyakran további fi\-nom\-han\-go\-lást igényel a végleges felhasználás előtt.


\vspace{2em}
Az első tézispont szempontjából releváns publikációk:
\begin{itemize}
	\item \textbf{Zoltán Ságodi}, István Siket, and Rudolf Ferenc
	\newblock {Methodology for Code Synthesis Evaluation of LLMs Presented by a Case Study of ChatGPT and Copilot}.
	\newblock \emph{IEEE Access}, VOL(12), 72303-72316, 2024.
\end{itemize}

\subsection*{A szerző hozzájárulásai}

%\input{Chapters/thesis_one_contrib_list.tex}
%%%%%%%% thesis_one_contrib_list


\begin{enumerate}[wide = 0pt, widest = {I/4.}, leftmargin =*]
	
	\item[I/1.] Összegyűjtöttem a kapcsolódó módszertanokat, és azonosítottam azok hiányosságait.
	
	\item[I/2.] Kidolgoztam az összehasonlítási módszertant.
	
	\item[I/3.] Megfelelő értékelési benchmarkot kerestem.
	
	\item[I/4.] Kiértékeltem a nyelvi modelleket a kiválasztott benchmarkon.
	
	\item[I/5.] Elvégeztem a szintetizált programok minőségi elemzését.
	
	\item[I/6.] Megterveztem a humán értékelést.
	
	\item[I/7.] Kidolgoztam a humán értékelési keretrendszert.
	
\end{enumerate}



%%%%%% theis one contrib list end


\section{\thesispointtwohu}

%Investigating hyperparameters, like the temperature of LLMs, is crucial, although it requires tremendous amounts of model inferences.
%Reusing model inferences provides a solution, although papers that perform model inference, often do not share the generated outputs~\cite{data_no_share_1,data_no_share_2,data_no_share_3,data_no_share_4,data_no_share_5}.
%We created a dataset that contains 18,900 raw and 18,896 processed LLM outputs in C++ source code generation.
%Infering LLMs, especially larger ones, requires special hardware which might not be available to everyone who wants to research LLMs.
%To broaden the possibilities of the research community, we evaluated 9 LLMs from 3 different LLM-families -- namely Llama, DeepSeek, and Qwen -- with various model sizes, including models with the size of 70 billion parameters.
%These models were inferred with variadic temperature settings ranging from 0.01 to 1.00.
A hiperparaméterek, például az LLM-ek hőmérsékletének, vizsgálata kulcsfontosságú, u\-gyan\-ak\-kor rendkívül nagy számú modell futtatást igényel.
A generált eredmények új\-ra\-fel\-hasz\-ná\-lá\-sa megoldást jelenthet, azonban az ilyen feladatokat végző publikációk gyakran nem osztják meg a generált kimeneteket~\cite{data_no_share_1,data_no_share_2,data_no_share_3,data_no_share_4,data_no_share_5}.
Ennek kezelésére lét\-re\-hoz\-tunk egy adathalmazt, amely 18.900 nyers és 18.896 feldolgozott LLM-kimenetet tartalmaz C++ forráskód-generálás területén.
Az LLM-ek futtatása avagy másnéven inferenciája, különösen a nagyobb modellek esetében, speciális hardvert igényel, amely nem minden kutató számára elérhető.
A kutatói közösség lehetőségeit bővítve, 9 különböző LLM-et vizsgáltunk három modellcsaládból,- Llama, DeepSeek és Qwen-, eltérő mo\-dell\-mé\-re\-tek\-kel, beleértve a 70 milliárd paraméteres modelleket is.
A modellek inferenciáját különböző hőmérsékleti beállításokkal végeztük, 0,01 és 1,00 közötti tartományban.

%The creation of this dataset did not include randomly selected values.
%We meticulously researched available open-source models and selected the best ones according to the HumanEval~\cite{codex} and MBPP~\cite{mbpp} benchmarks.
%To consider both benchmarks' results, we used a multi-condition method to select the best models.
%As the models were selected with great attention, the tasks to be generated are also considered with great attention.
Az adatkészlet összeállítása során nem véletlenszerűen választottunk modelleket.
Rész\-le\-te\-sen feltérképeztük az elérhető nyílt forráskódú modelleket, és a HumanEval~\cite{codex} és MBPP~\cite{mbpp} benchmarkok alapján kiválasztottuk a legjobban teljesítőket.
A két benchmark eredményeinek együttes figyelembevételéhez egy többfeltételes kiválasztási módszert alkalmaztunk.
Ahogyan a modellek kiválasztása, úgy a generálandó feladatok kiválasztása is nagy körültekintéssel történt.

%We needed tasks that are  meaningful for the LLMs, meaning they are not trivial to understand and implement.
%At first, using sources such as LeetCode seemed like a great solution, although such popular collections of tasks are most probably included in the training sets of LLMs.
%Therefore, these tasks also had to be hidden from the major benchmarks of LLMs to reduce the risk of overfitting.
%To achieve this goal, we used tasks from Sapientia ECN\footnote{\url{https://ecn.ms.sapientia.ro/}}, a programming competition.
%This competition guarantees that the tasks are non-trivial, and since the tasks are hand-made by the organizers, the originality of the tasks is also guaranteed.
Olyan feladatokra volt szükségünk, amelyek érdemben kihívást jelentenek az LLM-ek számára, vagyis nem triviálisak sem megértés, sem implementáció szempontjából.
Elsőre olyan források, mint a LeetCode, megfelelőnek tűntek, azonban ezek a népszerű feladatgyűjtemények nagy valószínűséggel szerepelnek az LLM-ek tanítóadataiban.
Ezért olyan feladatokra volt szükség, amelyek a főbb benchmarkokból is hiányoznak, ezzel csökkentve a túlillesztés kockázatát.
Ennek érdekében a Sapientia ECN\footnote{\url{https://ecn.ms.sapientia.ro/}} programozási verseny feladatait használtuk.
Ez a verseny garantálja, hogy a feladatok nem triviálisak, és mivel azokat a szervezők egyedileg készítik, az eredetiségük is biztosított.

%These tasks were provided to the LLMs alongside our prompt, which we created considering good prompting practices such as role specification or few-shot learning.
%The prompting techniques used in this prompt were role specification to narrow down the context, structured output requests to make the output easier to process, explanation enforcement to achieve more semantically coherent output, and few-shot learning to show the model what input-output pairs should be achieved by the generated program.

Ezeket a feladatokat a saját promptunkkal együtt adtuk a modelleknek.
A prompt\-ot bevált promptolási gyakorlatok figyelembevételével alakítottunk ki.
A promptban alkalmazott technikák közé tartozott a szerepmegadás a kontextus szűkítésére, a strukturált kimenet kérése a könnyebb feldolgozhatóság érdekében, a magyarázat kikényszerítése a szemantikailag koherensebb válaszok eléréséhez, valamint a few-shot tanulás, amely példákon keresztül mutatja meg a kívánt bemenet-kimenet párokat.

%We also validated the generated outputs by compiling the contained source code.
%Examples that compiled successfully were taken as successful source code generations, although those that failed the compilation had to be checked.
%A compilation failure could be the result of a real programming mistake made by the model, or it could be our fault by not extracting the correct or whole code part from the model-generated output.
%We manually checked a statistically significant number of the 3,959 examples and found that there were no mistakes in our framework; only the LLMs failed to generate source code that compiles.
A generált kimeneteket a bennük található forráskód lefordításával is validáltuk.
Azok az esetek, amelyek sikeresen lefordultak, sikeres kódgenerálásnak minősültek, míg a for\-dí\-tá\-si hibával végződő eseteket külön meg kellett vizsgálni.
A fordítási hiba eredhetett a modell által elkövetett programozási hibából, de akár abból is, hogy a generált kimenetből nem a megfelelő vagy nem a teljes kódrészletet emeltük ki.
A 3.959 példa egy statisztikailag szignifikáns részét manuálisan ellenőriztük, és megállapítottuk, hogy a keretrendszerünk nem tartalmaz hibát; a fordítási problémák kizárólag az LLM-ek hibájából adódtak.

%Examples that could be compiled were also evaluated with the available testcases from the programming competition.
%The evaluations resulted in multiple outcomes.
%Considering the whole task set over all the models, there were 29,882 testcases.
%The results are summarized in \autoref{th4:table:outcomes}.

A sikeresen lefordított programokat a versenyhez tartozó tesztesetekkel ki is értékeltük.
Az értékelések többféle kimenetet eredményeztek.
A teljes feladathalmazt és az összes modellt figyelembe véve összesen 29.882 teszteset állt rendelkezésre.
Az eredményeket a \ref{th4:table:outcomes} táblázat foglalja össze.



\begin{table}[h]
	\centering
	\begin{tabular}{|l|c|}
		\hline
		\textbf{Outcome} & \textbf{Number}  \\ \hline
		Passed           & 6.900            \\ \hline
		Failed           & 18.755           \\ \hline
		Runtime Error    & 2.306            \\ \hline
		Timeout          & 1.921            \\ \hline
		\textbf{Sum}     & \textbf{29.882}  \\ \hline
	\end{tabular}
	\caption{Number of different outcomes during program executions on test cases.}
	\label{th4:table:outcomes}
\end{table}

%As evident from the distribution in Table~\ref{th4:table:outcomes}, a significant portion of executions resulted in failures, with a smaller yet non-negligible number of runtime errors and timeouts.
%These results enable robust downstream analysis and interpretation tailored to diverse research or evaluation objectives.

Ahogyan az a \ref{th4:table:outcomes} táblázat eloszlásából is látható, a futtatások jelentős része hibával végződött, emellett kisebb, de nem elhanyagolható arányban fordultak elő futásidejű hibák és időtúllépések is.
Ezek az eredmények lehetővé teszik a megbízható további elemzést és értelmezést, különböző kutatási vagy értékelési célokhoz igazítva.

\vspace{2em}
A második tézispont szempontjából releváns publikációk:
\begin{itemize}
	\item \textbf{Zoltán Ságodi}, István Kolláth, Péter Heged\H{u}s and Rudolf Ferenc
	\newblock {A Program Synthesis Dataset for LLM Temperature Analysis}.
	\newblock \emph{IEEE Access}, VOL(13), 184022-184029, 2025.
\end{itemize}

\subsection*{A szerző hozzájárulásai}

%\input{Chapters/thesis_two_contrib_list.tex}

%%%%%% thesis two contrib list

\begin{enumerate}[wide = 0pt, widest = {II/5.}, leftmargin =*]
	
	\item[II/1.] Megterveztem a tanulmányt.
	
	\item[II/2.] Kidolgoztam a modellkiválasztási keretrendszert.
	
	\item[II/3.] Kidolgoztam a modellértékelési keretrendszert.
	
	\item[II/4.] Kiértékeltem az adatokat a hőmérséklet paraméter hatásának figyelembevételével.
	
\end{enumerate}



%%%%% tzheiss two contrib list end


\section{\thesispointthreehu}

%We took an alternative form of source code generation and used it to generate fixed source code to perform so-called automatic program repair.
%We used GPT-4 to generate the source code in a way that maintains the original functionality without the hidden vulnerability inside.
%This task was evaluated by Prenner et al.~\cite{codex_bug_fix} and Sobania et al.~\cite{quix_bugs_abuser_2}, although QuixBugs contains synthetic examples.
A forráskód-generálást alternatív módon alkalmaztuk, így javított forráskódot állítottunk elő úgynevezett automatikus programjavítás (automatic program repair) céljából.
A GPT-4 modellt arra használtuk, hogy olyan módon generáljon kódot, amely megőrzi az eredeti funkcionalitást, ugyanakkor megszünteti a rejtett sérülékenységet.
Ezt a feladatot korábban Prenner és társai~\cite{codex_bug_fix}, valamint Sobania és társai~\cite{quix_bugs_abuser_2} is vizsgálták, azonban a használt benchmark, QuixBugs, szintetikus példákat tartalmaz.

%In our work, this task is evaluated on a real-life vulnerability benchmark, Vul4J~\cite{vul4j}.
%This benchmark contains a real-life vulnerability collection alongside tests that validate whether the plausible fix correctly removes the vulnerability or not.
%To prevent our results from being influenced by GPT-4 being familiar with the vulnerabilities collected in the benchmark, We also used the CVEFixes~\cite{bhandari2021cvefixes} benchmark, which had no validation tests, thus We could only check correctness by hand.
%This benchmark has vulnerabilities only after the training period of that particular GPT-4 version.

Munkánkban egy valós sérülékenységeket tartalmazó benchmarkot, a Vul4J-t~\cite{vul4j}, hasz\-nál\-tuk kiértékelésre.
Ez a benchmark valós sérülékenységeket tartalmaz, valamint olyan teszteket, amelyek ellenőrzik, hogy az adott javítás valóban megszünteti-e a sé\-rü\-lé\-keny\-sé\-get.
Annak érdekében, hogy az eredményeket ne torzítsa az, hogy a GPT-4 esetleg ismeri a benchmarkban szereplő példákat, a CVEFixes~\cite{bhandari2021cvefixes} adatbázist is felhasználtuk.
Ez a benchmark nem tartalmaz validációs teszteket, így a helyességet csak manuálisan tudtuk ellenőrizni.
Azonban, az adatbázis olyan sérülékenységeket tartalmaz, amelyek az adott GPT-4 verzió tanítási időszaka után keletkeztek.

%We created our prompt by considering multiple prompting strategies and selected the best-performing on a subset of Vul4J.
%With our prompt, we instructed GPT-4 to fix the source code and also provide a textual answer.
%This two-sided request allowed us to discover if GPT can be a useful guide even if it cannot provide a valid vulnerability fix.
%Textual answers can be used as automatic pull request reviews or patching aids for the developers.
A promptoláskor több különféle stratégia szerint alakítottunk ki promptokat, és a Vul4J benchmark egy részhalmazán kiértékeltük azokat.
A legjobban teljesítő prompt-ot alkalmaztuk a teljes adathalmazra.
A prompt segítségével a GPT-4-et arra utasítottuk, hogy javítsa ki a forráskódot, és egyben adjon szöveges magyarázatot is.
Ez a kétoldalú meg\-kö\-ze\-lí\-tés lehetővé tette annak vizsgálatát, hogy a modell akkor is hasznos lehet-e i\-rány\-mu\-ta\-tás\-ként, ha helyes javítást nem is képes előállítani.
A szöveges válaszok felhasználhatók például automatikus pull request értékelésként vagy javítási segédletként a fejlesztők szá\-má\-ra.

%In contrast to many, in our evaluation, we used Top-1 evaluation, meaning we considered the first solution the model provided.
%This allows us to be closer to real-life scenarios, where developers do not have time to consider 5 or 10 plausible patches.
%We ran the evaluation three times, averaging the results, but all the distinct runs are considered as Top-1 evaluations.
Sok más megközelítéssel ellentétben az értékelés során Top-1 módszert alkalmaztunk, vagyis minden esetben a modell elsőként adott válaszát vettük figyelembe.
Ez közelebb áll a valós fejlesztési helyzetekhez, ahol a fejlesztőknek általában nincs idejük több (például 5-10) lehetséges javítás közül választani.
Az értékelést háromszor futtattuk le, és az eredményeket átlagoltuk, ugyanakkor minden egyes futtatást külön-külön is Top-1 ér\-té\-ke\-lés\-ként kezeltünk.

%Our findings show that in real-world scenarios the GPT-4 model can fix on average 33.33\% of the vulnerabilities and can provide useful answers in around 50\% of the cases, of which, in 16.67\%, only the text output was useful.
%We also investigated whether the results are consistent or not.
%Across the three individual runs, 27.03\% of the tests passed in every run.
%This consistency, however, is not related to the vulnerability type, but rather to good coding practices.
Az eredmények azt mutatják, hogy valós környezetben a GPT-4 modell átlagosan az esetek 33,33\%-ában képes a sérülékenységek javítására, és az esetek körülbelül 50\%-ában ad hasznos választ; ezek közül 16,67\%-ban kizárólag a szöveges válasz bizonyult hasznosnak.
Vizsgáltuk az eredmények konzisztenciáját is.
A három futtatás során az esetek 27,03\%-ában minden alkalommal sikeres volt a javítás.
Ez a konzisztencia azonban nem a sérülékenység típusával, hanem inkább a jó programozási gyakorlatok jelenlétével mutat összefüggést.


\vspace{2em}

A harmadik tézispont szempontjából releváns publikációk:
\begin{itemize}
	\item \textbf{Zoltán Ságodi}, Gábor Antal, Bence Bogenfürst, Martin Isztin, Péter Heged\H{u}s and Rudolf Ferenc
	\newblock {Reality Check: Assessing GPT-4 in Fixing Real-World Software Vulnerabilities}.
	\newblock In \emph{Proceedings of the 28th International Conference on Evaluation and Assessment in Software Engineering (EASE '24)}, Association for Computing Machinery, 252–261, 2024.
\end{itemize}

\subsection*{A szerző hozzájárulásai}

%\input{Chapters/thesis_three_contrib_list.tex}
%%%% thesis_three_contrib_list

\begin{enumerate}[wide = 0pt, widest = {III/5.}, leftmargin =*]
	
	\item[III/1.] Felkutattam a prompt-tervezési (prompt engineering) technikákat.
	
	\item[III/2.] Megterveztem a promptok és a modellek értékelési módszertanát.
	
	\item[III/3.] Részt vettem a generált forráskód manuális kiértékelésében.
	
	\item[III/4.] Részt vettem a generált szöveges válaszok manuális kiértékelésében.
	
\end{enumerate}

%%%%%%%%%% thesis three contrib list end

\section*{Összefoglalás}

%Software engineering is a rapidly evolving field.
%Developers have demonstrated, and will continue to demonstrate, a persistent need for source code analysis, both in direct and indirect forms.
%Direct forms include tools that assess code quality or detect vulnerabilities, while indirect forms encompass applications that leverage such analyses internally, such as IDEs or automated program repair (APR) tools.
%Although some features of these tools seem straightforward and easy to implement, they often rely on heavy research behind the scenes.
%These analyses are not just considered analysis; they can be split into two major categories: static and dynamic analysis.
%In this work we take static analysis, which uses the source code solely and usually does not need a compiled version of the code.
A szoftverfejlesztés egy gyorsan fejlődő szakterület.
A fejlesztők részéről folyamatos igény mutatkozik a forráskód-elemzésre, mind közvetlen, mind közvetett formában.
A közvetlen formák közé tartoznak azok az eszközök, amelyek a kód minőségét értékelik vagy sé\-rü\-lé\-keny\-sé\-ge\-ket azonosítanak, míg a közvetett formák olyan alkalmazásokat foglalnak magukban, amelyek belső működésük során támaszkodnak ezekre az elemzésekre, például az IDE-k vagy az automatikus programjavító (APR) eszközök.
Bár ezen eszközök egyes funkciói első ránézésre egyszerűnek tűnnek, a háttérben gyakran jelentős kutatási e\-red\-mé\-nyek állnak.
Ezek az elemzések két fő kategóriába sorolhatók: statikus és dinamikus elemzés.
Jelen dolgozat a statikus elemzésre fókuszál, amely kizárólag a forráskódot használja, és általában nem igényli annak lefordított változatát.

%Large Language Models can be considered a type of static analysis tool, as they can operate directly on source code without requiring compilation.
%LLMs are increasingly integrated into key software engineering processes, including design, documentation, and task management.
%Although these models are applied in diverse areas, this thesis focuses specifically on a form of source code synthesis.
A nagy nyelvi modellek (LLM-ek) a statikus elemzés egy speciális formájának te\-kint\-he\-tők, mivel képesek közvetlenül a forráskódon működni, nem igényelnek fordítást.
Az LLM-ek egyre inkább beépülnek a szoftverfejlesztés kulcsfolyamataiba, beleértve a tervezést, a dokumentációkészítést és a feladatkezelést.
Bár ezek a modellek számos területen alkalmazhatók, a jelen értekezés kifejezetten a forráskód-szintézis egy meghatározott formájára összpontosít.

%The aim of this thesis is to discover and deepen the connection between developers and LLMs as an ultimate development aid; therefore, this thesis is designed to discover many aspects of LLMs that could be useful for a developer, regarding possible comparison techniques.
%Among the many aspects, this thesis highlights the importance of hyperparameters, more precisely the temperature parameter, which is yet to be investigated thoroughly.
%To help such investigations, the thesis provides a dataset for researchers with multiple LLM-families and LLM-sizes. 
%This thesis also provides an actual status report, not just an abstract idea on what the realistic expectations are.

A dolgozat célja, hogy feltárja és elmélyítse a fejlesztők és az LLM-ek közötti kapcsolatot mint fejlesztést támogató eszköz.
Ennek érdekében a dolgozat több olyan szempontot vizsgál, amelyek hasznosak lehetnek a fejlesztők számára, különös tekintettel az összehasonlítási módszerekre.
A vizsgált tényezők közül kiemelt figyelmet kapnak a hiperparaméterek, ezen belül is különösen a temperature (hőmérsékelt) paraméter, amelynek hatása eddig még nem került kellő mélységben feltárásra.
Az ilyen jellegű vizsgálatok támogatására a dolgozat egy több LLM-családot és különböző modellméreteket tartalmazó adatkészletet is bemutat.
Emellett a dolgozat egy aktuális helyzetjelentést is nyújt, nem csupán elméleti megközelítést adva arról, hogy milyen reális elvárások fogalmazhatók meg ezekkel a modellekkel szemben.

\vspace*{3em}

%The author states that while the thesis results are primarily his own work, the pronoun We is used instead of I to recognize the input of co-authors in the papers forming the basis of this thesis.

A szerző kijelenti, hogy bár a dolgozat eredményei elsősorban saját munkáján alapulnak, az egyes szám első személy helyett többes szám első személyt használ (''mi''), ezzel is elismerve a dolgozat alapját képező publikációk társszerzőinek hozzájárulását.


% INFO: Get the publication list to a new page. If there is too much empty space left on the page before, then comment it out
\newpage

% Publication list
\undef{\chapter}
\newcommand{\chapter}{\section}
%\input{Chapters/my_publications}
%%%%%%% my publications

\chapter*{Publikációk}
\markboth{Publikációk}{Publikációk}
\addcontentsline{toc}{chapter}{Publikációk}

\vspace*{2em}

\subsection*{Folyóirat publikációk}

\vspace*{1em}

\begin{itemize}[wide = 0pt, widest = {[4]}, leftmargin =*]
	
	
	\item  \textbf{Zoltán Ságodi}, István Kolláth, Péter Heged\H{u}s and Rudolf Ferenc
	\newblock {A Program Synthesis Dataset for LLM Temperature Analysis}.
	\newblock \emph{IEEE Access}, VOL(13), 184022-184029, 2025.
	
	\item  \textbf{Zoltán Ságodi}, István Siket, and Rudolf Ferenc
	\newblock {Methodology for Code Synthesis Evaluation of LLMs Presented by a Case Study of ChatGPT and Copilot}.
	\newblock \emph{IEEE Access}, VOL(12), 72303-72316, 2024.
	
	\item  \textbf{Zoltán Ságodi}, Edit Pengő, Judit Juhász, István Siket and Rudolf Ferenc
	\newblock {Static Call Graph Combination to Simulate Dynamic Call Graph Behavior}.
	\newblock \emph{IEEE Access}, VOL(10), 131829-131840, 2022.
	
	\item  Edit Pengő, \textbf{Zoltán Ságodi}
	\newblock {A Preparation Guide for Java Call Graph Comparison}.
	\newblock \emph{Acta Cybernetica}, VOL(24), 131–155, 2019.
	
\end{itemize}

\subsection*{Konferencia publikációk}
\vspace*{1em}

\begin{itemize}[wide = 0pt, widest = {[4]}, leftmargin =*]
	
	\item  \textbf{Zoltán Ságodi}, Gábor Antal, Bence Bogenfürst, Martin Isztin, Péter Heged\H{u}s and Rudolf Ferenc
	\newblock {Reality Check: Assessing GPT-4 in Fixing Real-World Software Vulnerabilities}.
	\newblock In \emph{Proceedings of the 28th International Conference on Evaluation and Assessment in Software Engineering (EASE '24)}, Association for Computing Machinery, 252–261, 2024.
	
	\item  Judit Jász, István Siket, Edit Pengő, \textbf{Zoltán Ságodi}, Rudolf Ferenc
	\newblock {Systematic Comparison of Six Open-source Java Call Graph Construction Tools}.
	\newblock In \emph{Proceedings of the 14th International Conference on Software Technologies - ICSOFT}, SciTePress, pages 117-128, 2019.
	
\end{itemize}



%%%%%%%%%%%%% my publications end

% Make the numbering start from the end of the manual publication list
%\makeatletter
%\apptocmd{\thebibliography}{\global\c@NAT@ctr 10\relax}{}{}
%\makeatother
\newpage
\bibliographystyle{plain}
\bibliography{refs}


\end{document}

